% Source pins (SHA-256): PifoStatement.lean=fe9475ebfa2c462d27cb0a99781926074e1281396f2c9e8d32d6dfca788d814e; pifo-verify-all/VERIFICATION.md=c97ed5daa11c97f71a5d0e3867f96ddabe74438f1f598c049050c735a96b1887 (artifact audit; the universal calculation is proved below).
\section{A stateful separation}
\label{sec:stateful}

The fixed type-to-path assignment in \autoref{thm:main} is essential.  We now
work through two flat three-leaf trees whose assignment may inspect the
persistent arrival counter.  They agree on every flush word but disagree on a
five-operation interleaving.

For this example only, let a transaction for push number \(j\) choose a triple
\[
  \mathsf{tx}(j)=(\text{leaf},\text{leaf rank},\text{root rank}).
\]
The choice is independent of the packet label but depends on \(j\).  A pop does
not reset the arrival counter.  In \autoref{tab:stateful-assignments}, notation
\(C(1,1)\) means leaf \(C\), leaf rank \(1\), root rank \(1\).

\begin{table}[H]
\centering
\begin{tabular}{@{}c@{\qquad}l@{\qquad}l@{}}
\toprule
push number & scheduler \(S_1\) & scheduler \(S_2\)\\
\midrule
\(1\) & \(C(1,1)\) & \(E(1,3)\)\\
\(2\) & \(B(1,3)\) & \(E(2,1)\)\\
\(3\) & \(C(2,2)\) & \(F(1,2)\)\\
\(j\ge4\) & \(G(j,100+j)\) & \(G(j,100+j)\)\\
\bottomrule
\end{tabular}
\caption{Arrival-dependent assignments for the counterexample.  Leaf names
are local to a scheduler; both trees have three leaves.}
\label{tab:stateful-assignments}
\end{table}

\begin{proposition}[Stateful flush/interleaving separation]
\label{prop:stateful-separation}
The schedulers in \autoref{tab:stateful-assignments} have identical flush output
for every finite packet word.  They are not interleaved-equivalent.
\end{proposition}

\begin{proof}
First label packet occurrences by their source arrival numbers.  After pushes
\(1,2,3\), scheduler \(S_1\)'s root priority order is
\[
  C_1,\ C_3,\ B_2,
\]
because the respective root ranks are \(1,2,3\).  Leaf \(C\) contains
occurrences \(1,3\) at leaf ranks \(1,2\), while \(B\) contains occurrence
\(2\).  The drain is therefore \(1,3,2\).

Scheduler \(S_2\)'s root priority order is instead
\[
  E_2,\ F_3,\ E_1.
\]
Leaf \(E\) contains occurrence \(1\) at rank \(1\) and occurrence \(2\) at
rank \(2\); \(F\) contains occurrence \(3\).  The first root reference,
\(E_2\), selects leaf \(E\), but that leaf emits occurrence \(1\), not the
occurrence that created the root reference.  Then \(F\) emits \(3\), and the
last \(E\)-service emits \(2\).  The drain is again \(1,3,2\).  This is the
indirection visible in \autoref{fig:counterexample-hijack}.

For \(j\ge4\), both schedulers insert into the common tail leaf \(G\), with
root ranks \(100+j\) and leaf ranks \(j\).  These entries follow the first
three and are ordered increasingly, so both append \(4,5,\ldots,n\).  Directly
checking the shorter prefixes gives
\begin{equation}
 \begin{array}{c|c}
 n & \text{identity drain}\ \\
 \hline
 0 & \epsilon\\
 1 & 1\\
 2 & 1,2\\
 n\ge3 & 1,3,2,4,5,\ldots,n.
 \end{array}
 \label{eq:stateful-flush-pattern}
\end{equation}
Now take an arbitrary label word \(x_1\cdots x_n\).  The transactions ignore
labels, so replacing each identity in \autoref{eq:stateful-flush-pattern} by its
source label shows that both schedulers return
\(x_1x_3x_2x_4\cdots x_n\), with the displayed short cases.  Thus they agree
on every flush word, not merely on words with distinct labels.

For the interleaving
\begin{equation}
  \push;\ \popop;\ \push;\ \push;\ \popop,
  \label{eq:diverging-ops}
\end{equation}
the first push/pop emits occurrence \(1\) and empties both trees, but leaves
the counter at one.  Scheduler \(S_1\) then inserts occurrence \(2\) into
\(B\) at root rank \(3\) and occurrence \(3\) into \(C\) at root rank \(2\),
so the final pop emits \(3\).  Scheduler \(S_2\) inserts occurrence \(2\) into
\(E\) at root rank \(1\) and occurrence \(3\) into \(F\) at root rank \(2\).
The old rank-\(1\) packet in \(E\) is gone; the selected \(E\) leaf therefore
emits occurrence \(2\).  The identity traces are \(1,3\) and \(1,2\).

This is observable with a two-type alphabet.  Choose labels
\(x_1=0,x_2=0,x_3=1\).  The operation word
\(\push(0);\popop;\push(0);\push(1);\popop\) produces type traces
\(0,1\) for \(S_1\) and \(0,0\) for \(S_2\).  Hence interleaved equivalence
fails.
\end{proof}

\begin{figure}[t]
\centering
\begin{tikzpicture}[
  q/.style={draw,rounded corners,minimum width=3.25cm,minimum height=.72cm,
            align=center,font=\small},
  leaf/.style={draw,minimum width=1.7cm,minimum height=.62cm,font=\small},
  arr/.style={-{Stealth[length=2mm]}},
  note/.style={font=\scriptsize,align=center}]
  \node[q] (r1) at (-3.5,1.5) {$S_1$ root: $C_1,C_3,B_2$};
  \node[leaf] (c) at (-4.7,0) {$C:[1,3]$};
  \node[leaf] (b) at (-2.3,0) {$B:[2]$};
  \draw[arr] (r1) -- (c);
  \draw[arr] (r1) -- (b);
  \node[note,below=.22cm of c] {services $C,C$\\emit $1,3$};
  \node[q] (r2) at (3.5,1.5) {$S_2$ root: $E_2,F_3,E_1$};
  \node[leaf] (e) at (2.3,0) {$E:[1,2]$};
  \node[leaf] (f) at (4.7,0) {$F:[3]$};
  \draw[arr,very thick,pifored] (r2) --
    node[left,note,text=pifored] {$E_2$ selects $E$} (e);
  \draw[arr] (r2) -- (f);
  \node[note,below=.22cm of e,text=pifored] {first $E$ service\\emits $1$, not $2$};
\end{tikzpicture}
\caption{The three-push flush.  A root entry selects a child queue, not its
source packet.  The \(E_2\) reference can therefore release occurrence \(1\),
allowing the two different trees to realize the same permutation.}
\label{fig:counterexample-hijack}
\end{figure}

The artifact's executable model stores packet identity as the arrival number,
which is a stronger observation than type alone.  Lean kernel-checks the
separating interleaving and the flush equations through \(n=8\) in
\lean{Pifo.lean} using \lean{decide}; a separate \lean{native\_decide}
theorem in \lean{Extended.lean} checks them through \(n=30\).  The universal
flush claim in
\propref{prop:stateful-separation} is established by the symbolic calculation
\autoref{eq:stateful-flush-pattern}; the artifact explicitly does not contain an
all-\(n\) induction.  This boundary is recorded again in
\autoref{sec:mechanization}.

The standalone checker records only successful pop outputs, whereas the frozen
model records a failed pop as \(\None\).  Instrumenting it with observable
failures embeds the example in the paper's observation convention: every pop
in every flush and in the separating trace succeeds, so none of the displayed
outputs changes.
